\problema{Binary Tree Right Side View}{199}{src/02-arboles/199-right-side-view.cpp}{M}

Dado un árbol binario, imagina que te paras del lado \emph{derecho} del
árbol y miras hacia él. Queremos los valores de los nodos que alcanzas a
ver, de arriba hacia abajo (uno por cada nivel: el nodo más a la derecha de
ese nivel).

\paragraph{Intuición.} Si recorremos el árbol nivel por nivel (como en el
BFS de 994-Rotting-Oranges) y en cada nivel leemos los nodos de izquierda a
derecha, el nodo que se ve desde afuera es simplemente el \emph{último} que
leemos en ese nivel. No importa si ese nodo cuelga de la rama izquierda o
derecha de su padre; lo único que importa es su posición dentro del nivel:
el último en orden izquierda-a-derecha es, por definición, el más a la
derecha, y por lo tanto el visible.

\begin{center}
\ttfamily
\begin{tabular}{ccccc}
 & & 1 & & \\
 & 2 & & 3 & \\
 & & 5 & & 4\\
\end{tabular}
\quad$\rightarrow$\quad vista: [1, 3, 4]
\end{center}

\noindent Nótese el nivel 2: el nodo 5 (que cuelga de 2) queda oculto
detrás del nodo 4 (que cuelga de 3), simplemente porque 4 está más a la
derecha dentro de ese nivel.

\paragraph{Caso importante.} Un nivel puede tener su nodo visible colgando
de una rama \emph{izquierda}, si la rama derecha de ese nivel ya se quedó
sin hijos (es más corta):

\begin{center}
\ttfamily
\begin{tabular}{ccc}
 1 & & \\
\end{tabular}
\quad
\begin{tabular}{ccccc}
 & & 1 & & \\
 & 2 & & 3 & \\
 4 & & & & \\
\end{tabular}
\quad$\rightarrow$\quad vista: [1, 3, 4]
\end{center}

\noindent Aquí 4 cuelga de la izquierda de 2, pero en el nivel 2 no existe
ningún nodo del lado derecho: 4 es el único nodo de ese nivel y, por lo
tanto, el último leído de izquierda a derecha. Sí es visible desde la
derecha aunque provenga de una rama izquierda.

\paragraph{Solución.} BFS por niveles: en cada nivel guardamos el valor del
último nodo procesado (el de más a la derecha). También es posible resolverlo
con DFS visitando primero la rama derecha y luego la izquierda, registrando
el primer nodo encontrado en cada profundidad (el primero en ese orden de
visita es el más a la derecha; si ya hay un valor registrado para esa
profundidad, el nodo actual queda oculto detrás de uno visitado antes).

\lstinputlisting{src/02-arboles/199-right-side-view.cpp}

\complejidad{$O(n)$}{$O(n)$}

\paragraph{Notas de entrevista (Amazon).} El error típico es asumir que el
nodo visible de cada nivel siempre cuelga de una rama derecha; como muestra
el segundo diagrama, puede colgar de una izquierda si esa rama derecha
terminó antes. Aclarar esto en la entrevista demuestra que entendiste el
problema y no memorizaste una solución superficial. Casos borde a mencionar:
árbol vacío (vista vacía), árbol de un solo nodo (la vista es ese único
valor), y árboles muy desbalanceados donde toda una rama es más profunda que
la otra.
